\section{Separation of two global solutions}\label{sec:separation}

The preceding sections give two different ways to take limits. At a
fixed physical horizon, the finite solutions have global entropy limits
with local conservation. At a shrinking horizon, magnification at a
cap event produces a nonzero linear solution, and the branching argument
bounds its time interval. We now show that these facts force a cap limit
to differ from a dormant limit. The separation argument follows
\cite[Section~10]{OpenAIPeriodic}; the normalization here also removes
the factor \(m_T=T^p\).

All construction parameters are henceforth fixed, with the choices in
Proposition~\ref{prop:finite-dilated-horizon} and the subsequent smallness
requirements of the global estimates. Choose one dormant global limit
\(F^{\mathrm d}\) supplied by Proposition~\ref{prop:global-limits}.

\begin{lemma}\label{sep:dormant-vanishing}
On the periodic lift centered at the origin,
\[
 \frac{F^{\mathrm d}(Tt,Tx,v)}{T m_T}\longrightarrow0
 \qquad(T\downarrow0)
\]
in local \(L^1\) on
\[
 \mathcal V=(0,1)\times\R^3_x\times
                 \{v\in\R^3:\bar v\ne0\}.
\]
\end{lemma}

\begin{proof}
The early upper bounds for the dormant approximations are independent
of the approximation index. Their sum is locally integrable, so testing
against nonnegative functions passes these bounds to \(F^{\mathrm d}\).
We estimate them on a compact set on which
\[
 0<\tau_0\le t\le A<1,\qquad |x|+|v|\le R,
                  \qquad |\bar v|\ge\varepsilon>0.
\]
By \eqref{early:dormant-bound}, the hot upper bound is
\(C(Tt)^M H(v)\), for a fixed \(M>p+1\). Division by
\(T m_T=T^{p+1}\) makes it tend to zero uniformly relative to
\(H\) on this compact set.

The bath bound \eqref{early:bath-bound} is exponentially small here.
Indeed \(Tx\) lies inside the fixed initial hole with
\(D_0(Tx)\ge\delta\) for all sufficiently small \(T\). Hence its
normalized contribution is bounded by
\[
 C T^{-p-1}e^{-\beta'|v|^2}
       \exp\left(-\frac{c\delta^2}{TA}\right),
\]
which tends to zero.

For the cold upper bound, a main spatial enlargement at transverse
position \(Ty\), with \(y\) bounded, has scale \(s\le C_R T\).
For small \(T\), only the copy centered at the origin can be main
on this compact set: any other such copy would have a lattice center
with transverse distance at most \(s_0+o(1)\) from the origin and
axial distance at most \(a+b_{\rm sp}+o(1)<1/2\), contradicting
the unit minimum length of a nonzero lattice vector. The other copies
are covered by the nonmain bound.

The transverse Gaussian in the main-color bound of
Proposition~\ref{cold:control}, together with \(|\bar v|\ge\varepsilon\),
gives
\[
 c_{\rm main}(Tt,Tx,v)
 \le C s^{-31}
       \exp\left(-\frac{c\varepsilon^2}{d_0^2s^{20}}\right).
\]
The supremum of the right side over \(0<s\le C_R T\) tends to
zero faster than every power of \(T\). At each point there is at
most one main color. Lemma~\ref{cold:nonmain} bounds the remaining
sum by \(C_M(Tt)^M e^{-4|v|^2}\) for every fixed \(M>0\).
Taking \(M>p+1\) handles the same normalization for this sum.
All three components therefore vanish in local \(L^1\), as claimed.
\end{proof}

\begin{proof}[Proof of Theorem~\ref{thm:main}]
Take the datum \(F_0\) defined in \eqref{con:datum-equation}.
Lemma~\ref{con:datum} gives its finite mass, energy, and absolute
entropy integral, as well as bounded velocity support. For every
\(0<T\le T_0\), Proposition~\ref{early:cap-selection} gives cap
approximations, and Proposition~\ref{prop:global-limits} gives a global
entropy limit satisfying all the required local conservation laws and
having datum \(F_0\). The dormant limit \(F^{\mathrm d}\) has these
same properties.

Suppose there are horizons \(T_k\downarrow0\) at which a cap limit
equals \(F^{\mathrm d}\) almost everywhere on \((0,T_k)\).
For each fixed \(T_k\), choose approximations along the subsequence
converging to that particular common limit. Let \(K_j\) be a compact
exhaustion of \(\mathcal V\), and choose nonnegative smooth compact
functions \(\zeta_j\) on \(\mathcal V\) with
\(\zeta_j\ge1\) on \(K_j\). A diagonal choice of indices \(N_k\)
can ensure both \(s_{N_k}/T_k\to0\) and
\begin{equation}\label{sep:diagonal}
 f_k(t,x,v):=\frac{F_k(T_kt,T_kx,v)}{T_km_{T_k}}
       \longrightarrow0
       \quad\hbox{in local }L^1(\mathcal V),
\end{equation}
where \(F_k\) is the chosen finite solution.

Here is the diagonal choice explicitly. At stage \(k\), require
\(s_{N_k}/T_k<1/k\) and make the differences between the integrals
of the normalized approximation and normalized common limit against
\(\zeta_1,\ldots,\zeta_k\) smaller than \(1/k\). Each test, when
written in physical coordinates at the fixed horizon \(T_k\), is
a fixed smooth test of finite norm; periodizing its spatial factor
gives a test on the torus. Thus fixed-horizon weak convergence permits
these finitely many requirements. Lemma~\ref{sep:dormant-vanishing}
makes the integrals of the common limit tend to zero. Since the
approximations are nonnegative and \(\zeta_j\ge1\) on \(K_j\),
their integrals over \(K_j\) tend to zero. This proves
\eqref{sep:diagonal} without a convergence rate uniform in \(T_k\).

Choose a cap time \(t_{*,k}\) of each selected solution and set
\(r_k=t_{*,k}+s_{N_k}\). Proposition~\ref{prop:blowup-limit} gives,
after subselection,
\[
 \frac{s_{N_k}}{r_k}\to l,\qquad
 \frac{T_k}{r_k}\to U\in[1,\infty],\qquad
 \frac{h_k(r_kt,r_kx,v)}{m_{T_k}r_k}\rightharpoonup g
\]
locally in spacetime on \(0<t<U\). By
Proposition~\ref{prop:nonzero-limit}, \(g\not\equiv0\); by
Proposition~\ref{prop:finite-dilated-horizon}, \(U<\infty\).
Consequently
\[
 \rho_k=\frac{r_k}{T_k}\longrightarrow\rho=\frac1U>0.
\]
This positive lower bound permits comparison of the two dilations.
Since \(F_k=c_k+u_k+h_k\ge h_k\) through its cap horizon,
\begin{equation}\label{sep:scale-comparison}
 0\le\frac{h_k(r_kt,r_kx,v)}{m_{T_k}r_k}
       \le\rho_k^{-1}f_k(\rho_kt,\rho_kx,v).
\end{equation}

Fix a compact subset \(K\) of
\((0,U)\times\R^3_x\times\{\bar v\ne0\}\).
For all sufficiently large \(k\), the images of \(K\) under
\((t,x,v)\mapsto(\rho_kt,\rho_kx,v)\) lie in one fixed compact
subset \(K'\) of \(\mathcal V\). The four-dimensional
space--time Jacobian is \(\rho_k^4\), so
\eqref{sep:scale-comparison} yields
\[
 \int_K\frac{h_k(r_kt,r_kx,v)}{m_{T_k}r_k}\dd t\dd x\dd v
 \le\rho_k^{-5}\int_{K'}f_k(t,x,v)\dd t\dd x\dd v
 \longrightarrow0.
\]
Thus \(g=0\) off the velocity axis. The cap
\(0\le g\le(t+l)H\) makes \(g\) a density, and the axis has
Lebesgue measure zero. Hence \(g=0\) almost everywhere, contradicting
Proposition~\ref{prop:nonzero-limit}.

It follows that for all sufficiently small horizons no cap limit can
equal the fixed dormant limit on the corresponding early interval.
Choose one such cap limit. Both it and \(F^{\mathrm d}\) are global
entropy solutions with the same datum and all local conservation laws,
and they differ on a set of positive spacetime--velocity measure.
Rotating back from the lattice \(\Lambda\) to \(\mathbb Z^3\)
preserves these properties and gives the assertion on \(\mathbb T^3\).
\end{proof}
